% Rednote 海报的唯一排版母版；文档类由各章 poster-xx.tex 声明。
\input{../preamble}

\begin{document}
\newcommand{\PosterCompileDate}{%
    \number\year.\ifnum\month<10 0\fi\number\month.\ifnum\day<10 0\fi\number\day}
% 两行标题分别使用各章字号，保持明确换行且绝不自动折行。
\newcommand{\PosterTitleBlock}{%
    \vbox{\offinterlineskip
        \hbox{{\sffamily\bfseries
            \fontsize{\RednotePosterFirstLineSize}{\RednotePosterFirstLineSize}\selectfont
            \MakeUppercase{\RednotePosterFirstLine}}}%
        \vskip\PosterTitleLineGap
        \hbox{{\sffamily\bfseries
            \fontsize{\RednotePosterSecondLineSize}{\RednotePosterSecondLineSize}\selectfont
            \MakeUppercase{\RednotePosterSecondLine}}}%
    }%
}
\begin{tikzpicture}[remember picture,overlay]
    \fill[PosterAccent] (current page.south west) rectangle (current page.north east);
    % 花纹位于全部内容下方；白框随后覆盖花纹，保持信息区清晰。
    \begin{scope}
        \clip (current page.south west) rectangle (current page.north east);
        \node[inner sep=0pt,opacity=\PosterFlowerOpacity,rotate=\PosterFlowerUpperRotation]
            at ([xshift=\PosterFlowerUpperX,yshift=\PosterFlowerUpperY]current page.north east)
            {\includegraphics[width=\PosterFlowerUpperWidth]{../../_assets/images/flower-white.pdf}};
        \node[inner sep=0pt,opacity=\PosterFlowerOpacity,rotate=\PosterFlowerSideRotation]
            at ([xshift=\PosterFlowerSideX,yshift=\PosterFlowerSideY]current page.north west)
            {\includegraphics[width=\PosterFlowerSideWidth]{../../_assets/images/flower-white.pdf}};
        \node[inner sep=0pt,opacity=\PosterFlowerOpacity,rotate=\PosterFlowerLowerRotation]
            at ([xshift=\PosterFlowerLowerX,yshift=\PosterFlowerLowerY]current page.south east)
            {\includegraphics[width=\PosterFlowerLowerWidth]{../../_assets/images/flower-white.pdf}};
    \end{scope}

    % 顶部保留课程字标；GDP 柱图下移，与底部信息栏的黑色顶边合为一体。
    \begin{scope}[
        shift={([xshift=\dimexpr\PosterPageWidth-\PosterMargin-\PosterHeaderLogoWidth\relax,
                  yshift=-\PosterLogoTop]current page.north west)}
    ]
        \begin{scope}[scale=\PosterHeaderLogoScale,transform shape]
            \AdvMacroLogoText{PosterAccentText}{BodyText}
        \end{scope}
    \end{scope}

    \coordinate (title-origin) at
        ([xshift=\PosterTitleLeft,
          yshift=\PosterTitleVerticalShift]current page.west);
    \node[
        anchor=west,inner sep=0pt,text=BodyText,
        opacity=\PosterTitleShadowOpacity
    ] at ([xshift=\PosterTitleShadowX,yshift=\PosterTitleShadowY]title-origin)
        {\PosterTitleBlock};
    \node[
        name=title-main,anchor=west,inner sep=0pt,text=PosterAccentText,
        align=left
    ] at (title-origin) {\PosterTitleBlock};
    % 独立一行的序号直接跟随标题位置，字号自动等于第一行标题。
    \node[
        anchor=south west,inner sep=0pt,text=BodyText,
        font=\sffamily\bfseries\fontsize{\PosterChapterNumberSize}{\PosterChapterNumberSize}\selectfont
    ] at ([yshift=\PosterChapterNumberGap]title-main.north west)
        {\RednoteChapterDisplayNumber};

    % 底部信息栏：作者与日期共享容器和基线，各自保持清楚的阅读入口。
    \coordinate (footer-panel-sw) at
        ([xshift=\PosterMargin,yshift=\PosterFooterBottom]current page.south west);
    \coordinate (footer-panel-ne) at
        ([xshift=-\PosterMargin,yshift=\dimexpr\PosterFooterBottom+\PosterFooterPanelHeight\relax]current page.south east);
    \fill[PosterAccentText] (footer-panel-sw) rectangle (footer-panel-ne);
    % 柱图横轴与白框严格等宽，并直接充当白框顶边。
    \begin{scope}[
        shift={([yshift=\PosterFooterChartHeight]footer-panel-ne -| footer-panel-sw)}
    ]
        \begingroup
        \def\PosterLogoChartWidth{\PosterFooterChartWidth}%
        \def\PosterLogoChartHeight{\PosterFooterChartHeight}%
        \def\PosterLogoBarWidth{\PosterFooterChartBarWidth}%
        \def\PosterLogoAxisLineWidth{\PosterInfoRuleHeight}%
        \AdvMacroLogoPosterChart{BodyText}{BodyText}%
        \endgroup
    \end{scope}
    \node[
        anchor=north west,inner sep=0pt,text=PosterAccent,align=left,
        font=\sffamily\bfseries\fontsize{\PosterFooterMetaSize}{20}\selectfont
    ] at ([xshift=\PosterInfoPanelPadding,yshift=-\PosterFooterLabelInset]footer-panel-ne -| footer-panel-sw)
        {AUTHOR};
    \node[
        anchor=south west,inner sep=0pt,text=BodyText,align=left,
        font=\sffamily\bfseries\fontsize{\PosterAuthorSize}{36}\selectfont
    ] at ([xshift=\PosterInfoPanelPadding,yshift=\PosterFooterMainInset]footer-panel-sw)
        {Paul SUN};

    \node[
        anchor=north east,inner sep=0pt,text=PosterAccent,align=right,
        font=\sffamily\bfseries\fontsize{\PosterFooterMetaSize}{20}\selectfont
    ] at ([xshift=-\PosterInfoPanelPadding,yshift=-\PosterFooterLabelInset]footer-panel-ne)
        {COMPILED};
    \node[
        anchor=south east,inner sep=0pt,text=BodyText,align=right,
        font=\sffamily\bfseries\fontsize{\PosterDateSize}{35}\selectfont
    ] at ([xshift=-\PosterInfoPanelPadding,yshift=\PosterFooterMainInset]footer-panel-ne |- footer-panel-sw)
        {\PosterCompileDate};
\end{tikzpicture}
\null
\end{document}
